%% kapa: exam=1 round=25 year=2014 period=1 form=A subject=civil_law subject_ko=민법 seq=11 number=11
%% source: Q-Net 감정평가사 기출문제 (cst00309) artlSeq=5101207 file=extracted/1008/5101207/2014년 제25회 감정평가사 1차시험 기출문제/1교시(A형).hwp
%% conversion: kordoc 4.24.51 -> tools/gen1.py (auto); review=proofread-sonnet; flags=none
\begin{question}{11}
점유에 관한 설명으로 옳은 것은? (다툼이 있으면 판례에 의함)
\begin{choices}
  \item 무권리자로부터 토지를 매수한 자의 점유는 다른 특별한 사정이 없는 한 자주점유이다.
  \item 점유가 성립하기 위해서는 반드시 물건에 대한 물리적․현실적 지배가 수반되어야 한다.
  \item 건물의 소유자가 아닌 자가 실제로 건물을 점유하고 있다면 특별한 사정이 없는 한, 당연히 건물의 부지를 점유한다고 보아야 한다.
  \item 토지에 대한 소유권보존등기가 이루어졌다면, 그 등기명의자는 그 무렵 다른 사람으로부터 당해 토지에 대한 점유를 이전받았다고 본다.
  \item 부동산을 매도하고 등기를 이전하였으나, 아직 그 부동산을 인도하지 않은 매도인의 점유는 특별한 사정이 없는 한 자주점유로 본다.
\end{choices}
\end{question}
\answer{1}
